\documentclass[11pt,a4paper]{article}
\usepackage[margin=23mm]{geometry}
\usepackage{amsmath,amssymb,amsthm}
\setlength{\emergencystretch}{2em}
\newtheorem{lemma}{Lemma}
\newtheorem{theorem}[lemma]{Theorem}
\newcommand{\gen}[1]{\langle#1\rangle}
\newcommand{\Alt}{\operatorname{Alt}}
\title{Local pairs and ambient witnesses}
\author{Human Compression project: L044 continuation}
\date{}
\begin{document}\maketitle
Retain the actual odd-degree head setting: $G\le\Alt(\Omega)$
is transitive with connected string C-marking
$\rho_0,\ldots,\rho_{r-1}$, all maximal omissions
intransitive, $r\ge6$ and $n=2r+1$. Its first three
marks are an actual transitive head on $D=\{1,\ldots,6\}$,
and
\[
 \tau=\rho_3=(6\ b)\delta,\quad
 R=G_{\ge3}=\Alt(\Omega\setminus\{1,\ldots,5\}),\quad
 \rho_4|_D=1,\quad H=G_{\ge5},\quad H|_{D\cup\{b\}}=1.
\]
The odd involution $\delta$ is supported outside $D\cup\{b\}$.
Put $K=G_{>5}$. Assume every colour $j\ge5$ has two
eligible edges for its full ambient omission $G_j$,
without assuming they lie in the same $H$-orbit.
For an actual $H$-orbit $O$, its \emph{local omission}
of $j$ is the restriction of
$\gen{\rho_i:i\ge5,\ i\ne j}$ to $O$.

\begin{lemma}[The weaker assignment interface]
It is impossible to assign each colour $j\ge5$ to
one actual $H$-orbit $O_j$ so that all three of the
following conditions hold:
\begin{enumerate}
\item $O_j$ contains an ambient-eligible $j$-edge;
\item $O_j$ contains two $j$-edges eligible for its
      full local omission;
\item $K$ is intransitive on $O_j$.
\end{enumerate}
The two edges in the second condition need not be
ambient eligible.
\end{lemma}
\begin{proof}
We verify the input weakening in the complete acquired
odd-degree orbit-budget proofs. Their stronger
assignment hypotheses are not invoked as theorem
statements here. Instead we replay their proofs using
the following precise interfaces.

For every nontrivial $H$-orbit $O$, let $F_O$ be
the assigned colours and $d_O=|F_O|-|O|/2$.
The selected-colour bound applies to the transitive
action $H|_O$, using its full local eligible pairs
from condition 2. It gives $d_O\le0$ and makes
$G_{F_O}$ transitive on $O$ in equality. If $5\notin F_O$,
then $G_{F_O}\le K$ is intransitive by condition 3,
so equality is impossible. This recovers all the
normalized deficit inequalities, including their
integral and half-integral forms. Empty assignments
retain their exact deficit. The fixed-anchor proof
also survives verbatim: it uses this local equality
and commuting marked components, and does not need
an additional ambient edge.

Every internal pair forest chooses exactly one
ambient-eligible edge for each assigned colour.
Condition 1 supplies it. The pair switches are
provided by $\tau\in G_j$. Thus the complete
internal-pair theorem, its equality path and faithful
rows, and the odd-degree row exit retain their
stated hypotheses. For two pairs, the same direct
argument shows that an eligible edge exchanges the
pair vertices; a second tail mark exchanging them
would connect all four points inside $G_j$.
Consequently the two-row proof likewise needs only
condition 1 and actual $K$-intransitivity.

When $\delta$ exchanges $H$-orbits, transport the
two \emph{local} edges by the actual $\tau$.
It centralizes $H$ and each local omission, so
maps their orbit partitions bijectively between
the two orbits. Local eligibility is preserved.
Therefore the selected-colour bound applies to
the union of their assigned colours in one orbit,
with exactly the combined deficit used in the
acquired proofs. The ambient witness also remains
ambient, since $\tau\in G_j$.

All localization and terminal-interval counts use
two ambient eligible edges \emph{globally}, which
are a standing hypothesis here. In the acquired
proofs the edges of the relevant colours are shown
to be confined to a displayed union of orbits:
on the complementary internal orbit their pair
images are trivial, and every possible edge there
is already joined by $\tau$ in its ambient omission.
This confinement uses the ambient witness forest,
not the local pair from condition 2. Thus the
complete global localization bounds remain available.

For clarity, this accounts for every assignment use
in the remaining special cases. In the first internal
orbit proof, pair-path equality confines the remaining
ambient edges to the companion orbit; its second
fixed anchor and the literal two-point residual are
unchanged. In the nonfirst deficit-two proof, the
local equality on $V$ supplies transitive $G_{F_V}$
and the pure $\Alt(V)$ prefix kernel. One ambient
witness per colour supplies the two-component
pair forest. The kernel, faithful projection,
connected marking and three-point row-growth
arguments are therefore unchanged. Its final
four-point contradiction uses two distinct actual
edges of $\rho_6$ on $V$, which condition 2 supplies.
The earlier four-point-row terminal contradiction
uses the same fact. No eligibility beyond existence
of these two edges is used at that step.

The actual head overlap and pure-prefix kernels
depend on none of the assignment pairs. The literal
exterior-transposition obstruction requires two
ambient edges for each later colour globally,
already assumed. These observations exhaust all
uses of a second assigned ambient edge in the
eight-fixed-point proof, the first and nonfirst
internal proofs, the fixed-anchor corollary, and
the last fixed-point budgets.

We can now replay the exhaustive normalized assembly:
at seven fixed points, exchanges, the two first-orbit
budgets, the nonfirst deficit-two budget and the empty
internal pair are excluded by exactly those arguments.
The eight-point argument is unchanged as just checked.
Hence $|\operatorname{Fix}(H)|\ge9$. The last-budget
proof uses the global interval bound $f+c\le11$,
the actual anchor which is not $H$-fixed, the same
local deficits and internal witness forests. Its
sole remaining four-point row has a terminal label
by local fixed-anchor equality, and its orbit growth
is disjoint from ten actual fixed points. It gives
the opposite inequality $|\operatorname{Fix}(H)|<9$.
This contradiction proves the weaker-interface lemma.
\end{proof}

\begin{lemma}[Active labels at a single ambient edge]
Let $O$ be an actual $H$-orbit containing exactly one
ambient-eligible edge of some colour $j\ge5$.
The nonidentity restrictions of the marked generators
to $O$ have connected noncommutation diagram.
Consequently their active labels form a consecutive
interval containing $j$.
\end{lemma}
\begin{proof}
Split that diagram into its connected components.
The corresponding permutation groups commute and
generate the transitive action on $O$. Any component
not containing $j$ centralizes $\rho_j|_O$ and lies
in the actual ambient omission $G_j$. It sends the
unique eligible edge to itself: it preserves $O$,
commutes with the edge involution and preserves
each ambient omission orbit. Its endpoints lie in
different such orbits, so each is fixed individually.

A factor of a transitive commuting product which
fixes one point fixes every point. Indeed, remove
that factor from a word carrying the fixed point
to an arbitrary point; the other factors commute
with it and the removed factor fixes the initial
point. Every other component would therefore be
trivial on $O$, contradicting its definition.
Only the component containing $j$ remains.
Noncommutations in a string are possible only
between consecutive indices, proving the interval
assertion.
\end{proof}

\begin{lemma}[The first active label]
Suppose each actual $H$-orbit has at most one
ambient-eligible edge of colour $j$. On any orbit
containing such an edge, the active interval of
the preceding lemma starts at $5$ or $6$.
\end{lemma}
\begin{proof}
An orbit on which $K$ is intransitive has a nontrivial
$\rho_5$-restriction, and its interval starts at $5$.
Consider instead a $K$-transitive orbit $O$ whose
first active index $m$ is greater than $5$.
Suppose $m\ge7$. The whole-orbit transport supplies
an actual $K$-orbit $T=O^p$ inside a $K$-intransitive
$H$-orbit $W$, with $p\in G_{<5}$. It transports
the ambient $j$-edge and the complete marked
$K$-action. Thus $\rho_6$ fixes $T$ and
$L=G_{\ge7}$ is transitive on $T$.

The mark $\rho_5$ centralizes $L$, so $T^{\rho_5}$
is an actual $L$-orbit. If it differs from $T$,
they are disjoint subsets of $W$. Since $j\ge m\ge7$,
$\rho_5$ transports the witness to a second
ambient-eligible $j$-edge in $W$, a contradiction.
Hence $\rho_5$ preserves $T$. It preserves the unique
ambient $j$-edge in $W$ and fixes its endpoints
individually, since it belongs to $G_j$. Centralization
of the transitive $L|_T$ then makes it fix $T$
pointwise. Both $K$ and $\rho_5$ preserve $T$, so
$W=T$, contrary to the chosen $K$-intransitivity
of $W$. Thus $m=6$ is the only possibility.
\end{proof}

\begin{theorem}[A genuine local crossing after transport]
Some colour $j\ge5$ and actual $H$-orbit $O$ satisfy:
\begin{enumerate}
\item its full local $j$-omission has exactly two
      orbits, joined by exactly one $j$-edge;
\item that edge is ambient eligible, and every
      actual $H$-orbit has at most one ambient-eligible
      $j$-edge;
\item $K$ is intransitive on $O$, and the active
      labels on $O$ form an interval beginning at $5$.
\end{enumerate}
There is a second, distinct $H$-orbit with an
ambient-eligible $j$-edge. Its active labels form
an interval beginning at $5$ or $6$; local
uniqueness and $K$-intransitivity on that second
orbit are not asserted.
\end{theorem}
\begin{proof}
For each colour choose an $H$-orbit with two ambient
eligible edges if one exists, and retain both.
Otherwise choose any one ambient witness edge.
If $K$ is transitive on that orbit, the complete
whole-orbit transport gives an element of $G_{<5}$
carrying it, as a whole $K$-orbit, into an $H$-orbit
on which $K$ is intransitive. Transport every retained
edge by the same element. A colour needing transport
is greater than $5$: an ambient $5$-edge cannot
lie in a $K$-transitive orbit. The transporting
element commutes with the colour and lies in its
ambient omission, so all retained edges remain
actual and ambient eligible. The preferences are
therefore preserved simultaneously.

Each final assigned orbit contains at least one
edge eligible for the full local omission. If all
contained two, they would satisfy the preceding
weaker-interface lemma, which is impossible.
Choose one with exactly one. A transitive action
with exactly one edge joining different full
omission orbits has exactly two such orbits:
the connected graph obtained by collapsing them
has just that one edge. This proves statement 1.

If any $H$-orbit had two ambient $j$-edges, the
preference and simultaneous transport would have
supplied two in this target, hence two locally
eligible edges. This is impossible, proving 2.
The target was normalized to be $K$-intransitive.
Thus $\rho_5$ acts nontrivially on it; otherwise
$K$ and $H$ would have the same transitive action.
The previous lemma proves 3.

There are two ambient $j$-edges globally, and
statement 2 puts another one in a different
$H$-orbit. The preceding two lemmas give its
active interval and its first index. We do not
transport this second edge and claim its image
is distinct: such transported witnesses may
coincide, and no such assertion is needed.
\end{proof}

\paragraph{Attribution and boundary.}
The complete acquired case proofs remain bound to
their original hypotheses. The first lemma explicitly
checks each changed use and proves the weaker
interface by their complete arguments. The new
transport preference and commuting-factor arguments
then supply a genuine local unique crossing, rather
than inferring one from a split ambient pair.
The source's local crossing and two-path strategy
is retained, with these orbit distinctions made
explicit. Full path counting, terminal residuals
and later global attachments remain outside this
note. No new induction, computation, formal,
original or question credit is claimed.
\end{document}
